\section{Results and discussion} \label{sec:resultsNdiscussion}
\subsection{Econometric results}

\begin{figure}[tb] 
	\centering
	\includegraphics[width=\textwidth]{overall_rif_decompose}
	\caption{RIF decomposition}
	\label{fig:overallrif}
\end{figure} 
 
Table \ref{decomposemanager}, \ref{decomposeclerical} and \ref{decomposeelem} reports the estimates of RIF-OLS decomposition and presents the baseline result for managers, clerical workers and elementary workers respectively, with annual normalized wage earnings for 2008 and 2018 as the dependent variable. We report results for the total effect, unexplained and explained effects along with their estimates, standard errors and bounds. The results indicate that informal employees suffer an earnings penalty across  the unconditional earnings distributions for managers and clerical occupation for both the year of 2008 and 2018. For the elementary occupation, there is wage premium for the informal employees for upto 30th quantile whereas wage penalty above the 30th quantile for 2018 and wage penalty across the distribution for the year 2008. Specifically, the raw earnings gap is in the order of 0.427 and -0.132 for elementary occupation, 0.76 and 0.414 for clerical occupation and 0.759 and 0.59 for managers at 10th quantile  for 2008 and 2018 respectively.  Similarly, the raw earnings gap is in the order of 338 and 0.05 for elementary occupation, 0.815 and 0.562 for clerical occupation and 0.591 and 0.611 for managers at 50th quantile for 2008 and 2018 respectively. Moreover, the raw earnings gap is in the order of 0.354 and 0.209 for elementary occupation, 0.598 and 0.279 for clerical occupation and 0.207 and 0.207 for managers at 10th quantile  for 2008 and 2018 respectively. An observation, that can be drawn from these is that raw earnings gap is decrasing for all the occupations from year 2008 to 2018 across all the uncoditional distributions, except for 30-50th quantile for managers as shown in table \ref{decomposemanager}. 

We plot these results as well as the overall wage gap and composition and endowment effects in figure \ref{fig:overallrif}. The results shows that wage gap for elementary occupation is highest at the 10th quantile, starts decreasing and increases after the median in 2008 whereas it is highest at 80th quantile, negative below 30th quantile and minimal upto 60th quantile. For the clerical occupation, wage gap is highest at the 50th quantile, decreases above median, but decreases at an increasing rate below median and increases at the 10th quantile in 2008, whereas highest at 40th and 50th quantile and decreases below and above median for 2018. This is evidence of the presence of the sticky floor phenomenon in the informal labour market for elementary occupation but not for the rest of the occupations. These findings are similar to other studies in the developing economy. 

The unconditional quantile decomposition results further indicate that the 









   
